% §18.6 Closing the Book
%
% Monograph register: the five-part arc (atom; sets; maps and types;
% constructive optimum; boundary and applications) and the research
% directions. Tightened; the arc and chapter mapping are clean and
% preserved. Cross-refs the chapter labels and chapter-4 axiom/theorem
% labels. DEFINED: sec:closing-the-book.

\section{Closing the Book}
\label{sec:closing-the-book}

The book built one framework in five parts and nineteen chapters, each
part lifting the preceding part's algebra to a larger domain; chapter~19,
which follows, restates the whole in operator form. This section traces
the arc and says plainly what kind of contribution it is.

\subsection*{The atom and its lifts}

Part~I opened with the noisy bit (\Cref{ch:noisy-bit}), a Boolean value
crossing a binary channel, and named the two atomic parameters, the
false-positive rate $\fprate$ and the false-negative rate $\fnrate$,
jointly inhabiting the Bernoulli square: the asymmetric binary channel\index{channel!asymmetric binary} of
information theory and the randomized response\index{randomized response} of survey statistics.
\Cref{ch:axioms-kronecker} formalized the atom: $\Bernoulli{\BoolType}$
(\Cref{def:bernoulli-bool}) collects all $(\fprate, \fnrate)$ channels, the
two axioms (\Cref{ax:element-indep}, \Cref{ax:block-rates-indep}) prescribe
how atoms combine, and the Kronecker theorem (\Cref{thm:kronecker}) gives
the joint law of parallel queries as a tensor product. Serial stages
compose by the matrix product, and Boolean formulas compose exactly per
latent assignment, not through a single carried rate pair
(\Cref{ex:two-rate-counterexample}).

Part~II lifted the atom to a set of bits, one per element of a universe.
\Cref{ch:approximate-sets} framed the random approximate set as a function
from a plaintext set to a sequence of atoms; \Cref{ch:set-composition}
worked out the rates of union, intersection, complement, and difference,
with the perfect-filter line marking the structures (Bloom filters chief
among them) on the $\fnrate = 0$ axis;
\Cref{ch:classification-measures} lifted the atomic rates to the derived
measures (precision, recall, PPV, NPV, $\Fone$, Youden's $J$); and
\Cref{ch:information-bounds} recalled the space floor of approximate
membership, $\log_2(1/\fprate)$ bits per element, with its two-sided
extension.

Part~III lifted from sets to functions. \Cref{ch:approximate-functions}
took the binary channel to its $n$-ary generalization, the random
approximate map; \Cref{ch:sketches-estimators} read streaming sketches in
the same terms; \Cref{ch:randomized-algorithms} placed randomized
algorithms in the framework; \Cref{ch:approximate-relations} turned to
query-level error contracts for approximate relations; and
\Cref{ch:algebra-types} organized them as instances of $\Bernoulli{T}$
for types $T$, with the limits of that organization (shared intermediate
values, sum types) stated.

\subsection*{The constructions, the boundary, the applications}

Part~IV asked how close constructions come to the space floor.
\Cref{ch:optimal-bernoulli-sets} treated the FNR-zero arm with a perfect
hash addressing a fingerprint table, the construction
\textcite{dietzfelbinger-pagh-2008} describe, and seed search in the line
of RecSplit and ShockHash; \Cref{ch:perfect-hash-filters} treated the
FPR-zero arm and ranked retrieval. Part~V drew the boundary and turned the
algebra to use. \Cref{ch:not-bernoulli} catalogued where the framework
fails, when rates drift, errors correlate, inputs are adversarial, or no
latent value exists, and what replaces the broken machinery.
\Cref{ch:cipher-maps-trapdoor} read keyed approximate structures as cipher
maps, with the chapter 7 posterior as one measure of the cover false
positives provide, and \Cref{ch:encrypted-search}, the present chapter,
applied that reading to Bloom-filter encrypted search.

\subsection*{What kind of contribution this is}

Most individual results in these chapters have named predecessors, and
each chapter's Notes and Lineage section credits them: randomized
response and misclassification models for the atom, Carter et al.\ and
Pagh and Rodler for the space bounds, Parker and McCluskey and
probabilistic transfer matrices for gate-level propagation, Gallager for
the parity identity, Dietzfelbinger and Pagh for the perfect hash filter,
probabilistic databases for relational error, and Decaf and its relatives
for probability-carrying types. What the book adds is a single notation
and algebra that runs through all of them, a map of where that algebra
holds exactly, where it holds only as an average or a bound, and where it
fails, and an executable companion that checks the quantitative claims.

\subsection*{What comes next}

Three threads point past the current scope. Approximate query languages
could lift the chapter 12 error contracts to a declarative interface with
calibrated per-source budgets. Multi-party settings could extend cipher
maps to several parties holding partial keys. And the relation to
differential privacy deserves a careful statement: randomized response,
the atom of this book, is also the canonical local differential-privacy
mechanism \parencite{warner-1965,erlingsson-pihur-korolova-2014}, and
differential privacy constrains ratios of channel entries rather than the
entries themselves, so the two vocabularies describe the same channels
under different guarantees rather than one subsuming the other. Each is a
possible extension, not a planned chapter.
